\begin{tcolorbox}[colback=red!3!white,colframe=red!50!black,boxrule=0.8pt,title={Correction notes},fonttitle=\bfseries]
\textbf{Corrections from the calculation review (30~Sept.~2026).} The following places are corrected or annotated in the text; each carries a dated note there with the previous value:
\begin{itemize}
\item Characteristic energy $E_0$: gravitational-geometric derivation $E_0^2 = 4\sqrt2\,m_\mu/\xi^4$ ($1.4 \times 10^{9}$~MeV, dimensionally inconsistent) and the section \enquote{Two Geometric Paths to $E_0$} removed; $E_0 = \sqrt{m_e m_\mu/K_{\text{frac}}}$ (R72).
\item Strong CP phase (Level~7 table): $\theta_{QCD} = \xi^2 = 1.78 \times 10^{-8}$ lies about 178 times above the bound $10^{-10}$ and is excluded (note).
\item Precision correction of $E_0$: $7.35 \times (1 + \alpha/2\pi) = 7.358$~MeV $\to$ $7.348/\sqrt{K_{\text{frac}}} = 7.397$~MeV (R72); hence $\alpha^{-1} = 137.06$.
\item Derivation via $D_{\text{frac}} = 4.2 \times 10^{-5}$: the formula gives $\approx 1.15 \times 10^{20}$; the value is fitted to $137.036$; annotated.
\item Table level~1: $\alpha_S = \xi^{-1/3}$: $9.65 \to 19.57$; $\varepsilon_T = \xi\sqrt{3/(4\pi^2)} = 3.7 \times 10^{-5} \to \varepsilon_T = \xi E_0^2 = 7.297 \times 10^{-3}$ (also in note (*)).
\item Table: $\sin^2\theta_W$ from $\alpha_W = \sqrt{\xi}$: $0.231 \to 0.0058$; $\delta_{CP} = \pi(1 - 2\xi)$: $1.57 \to 3.14$ rad.
\item Higgs VEV appendix: derivation from \enquote{Derivation of Higgs VEV} onward removed (conversions $245.1$ and $141.0$~GeV wrong, $K_{\text{quantum}}$ fixed from the measured value); replaced by the structural formula $v = \xi^4E_P/(5\pi) = 245.65$~GeV (Docs.~149, 383), likewise in the table.
\item Tables levels~2--5: $\Lambda_{QCD} = v\xi^{1/3}$: $200$~MeV $\to$ $12.6$~GeV; $m_h = v\xi^{1/4}$: $125 \to 26.4$~GeV; $m_\tau$: $\frac54 \to \frac{25}{9}$, $1778 \to 1783$~MeV; $m_c$: $\frac89 \to 2$, $1.279 \to 1.284$~GeV; neutrinos $\sim10^{-3}/10^{-2}/10^{-1} \to \sim10^{-6}/10^{-4}/10^{-3}$~eV. $\kappa_{\text{mass}} = D_f/2 = 1.47 \to D_f^{\text{eff}}/2 = 1.487$. Direct and Yukawa methods not equivalent ($3.2$ vs.\ $207.8$), muon $l = 1$ vs.\ $l = 0$ (note).
\item (1~Oct.~2026) Table level~6 (K-041-a/b/c): $\delta_{CKM} = \arcsin(2\sqrt2\,\xi^{1/2}/3)$: $1.2 \to 0.0109$~rad; $\theta_{13} = \arcsin(\xi^{1/3})$: $8.6° \to 2.93°$; $|V_{us}|$: $0.225 \to 0.21$.
\item (1~Oct.~2026) Table level~4: $m_s = v \cdot 3 \cdot \xi$: $97.9 \to 98.4$~MeV (Doc.~006: $r_s = 26/9$, $94.76$~MeV).
\end{itemize}
\end{tcolorbox}

\hfuzz=10pt
\section*{Abstract}
		This documentation presents the derivation of the parameters of T0-theory in logical order, with the aim of excluding circularity. The presentation shows how the fine structure constant $\alpha = 1/137$ is meant to follow from geometric principles without presupposing it. The derivation steps are documented explicitly so that the charge of circularity can be checked step by step. According to the present state of the corpus, the early coupling-constant formulas of this document are classified as preliminary stages (see the addendum at the end, R77).

	\medskip\noindent\textbf{Note on versions.} The German and English versions of this document differ in structure and extent; earlier revisions shortened, extended or changed parts of both. Where they differ, the more extensive German version applies; sections contained only in this English version count as supplements.


\medskip\noindent\textbf{Status markers.} Statements marked \textbf{[S]} are \emph{speculative}: a hint or conjecture that has not yet been derived or numerically checked in the corpus.


	\section{Introduction}
	
	T0-theory takes the approach of treating fundamental physical constants not as arbitrary but as derivable from the geometric structure of three-dimensional space. The central claim is that the fine structure constant $\alpha = 1/137.036$ is not an empirical input but follows from spatial geometry.
	
	To make the question of circularity checkable, we present here the derivation of the parameters in logical sequence, starting from geometric principles and without using experimental values except fundamental natural constants; where a measured value does enter at some point, this is noted in the text.
	
\section{The Geometric Parameter $\xi$}

\subsection{Derivation from Fundamental Geometry}

The universal geometric parameter $\xi$ consists of two fundamental components:
\begin{equation}
	\xi = \frac{4}{3} \times 10^{-4}
\end{equation}

\subsubsection{The Harmonic-Geometric Component: 4/3 as the Universal Fourth}

\textbf{4:3 = the fourth -- a universal harmonic ratio}

The factor 4/3 is interpreted here as the \textbf{perfect fourth}, one of the basic harmonic intervals:

\begin{equation}
	\frac{4}{3} = \text{Frequency ratio of the perfect fourth}
\end{equation}

Just as musical intervals are universal:
\begin{itemize}
	\item \textbf{Octave:} 2:1 (always, whether string, air column, or membrane)
	\item \textbf{Fifth:} 3:2 (always)
	\item \textbf{Fourth:} 4:3 (always)
\end{itemize}

These ratios are \textbf{geometric/mathematical}, not material-dependent.

\textbf{Why is the fourth universal?}

For a vibrating sphere:
\begin{itemize}
	\item When divided into 4 equal ``vibration zones''
	\item Compared to 3 zones
	\item The ratio 4:3 emerges
\end{itemize}

This is \textbf{pure geometry}, independent of material.

\textbf{The harmonic ratios in the tetrahedron:}

The tetrahedron contains BOTH fundamental harmonic intervals:
\begin{itemize}
	\item \textbf{6 edges : 4 faces = 3:2} (the fifth)
	\item \textbf{4 vertices : 3 edges per vertex = 4:3} (the fourth)
\end{itemize}

\textbf{The complementary relationship:}
Fifth and fourth are complementary intervals - together they form the octave:
\begin{equation}
	\frac{3}{2} \times \frac{4}{3} = \frac{12}{6} = 2 \quad \text{(Octave)}
\end{equation}

The document reads this as a harmonic structure of space:
\begin{itemize}
	\item The tetrahedron contains both fundamental intervals
	\item The fourth (4:3) and fifth (3:2) are reciprocally complementary
	\item The harmonic structure is self-consistent
\end{itemize}

\textbf{Further appearances of the fourth in physics:}
\begin{itemize}
	\item Crystal lattices (4-fold symmetry)
	\item Spherical harmonics
	\item The sphere volume formula: $V = \frac{4\pi}{3}r^3$
\end{itemize}

\textbf{Interpretation:}
\begin{itemize}
	\item \textbf{Parallel to Pythagoras:} ``Everything is number and harmony''
	\item \textbf{Space itself} has a harmonic structure
	\item \textbf{Particles} are ``tones'' in this cosmic harmony
\end{itemize}

T0 theory thus interprets space as harmonically structured, with 4/3 (the fourth) as its basic signature. This interpretation motivates the choice of the factor 4/3; it is not a derivation in the strict sense.

\textbf{The $10^{-4}$ Factor:}

\textbf{Step-by-Step QFT Derivation:}

\textbf{1. Loop Suppression:}
\begin{equation}
	\frac{1}{16\pi^3} = 2.01 \times 10^{-3}
\end{equation}

\textbf{2. T0-Calculated Higgs Parameters:}
\begin{equation}
	(\lambda_h^{\text{(T0)}})^2 \frac{(v^{\text{(T0)}})^2}{(m_h^{\text{(T0)}})^2} = (0.129)^2 \times \frac{(246.2)^2}{(125.1)^2} = 0.0167 \times 3.88 = 0.0647
\end{equation}

\textbf{3. Missing Factor to $10^{-4}$:}
\begin{equation}
	\frac{10^{-4}}{2.01 \times 10^{-3}} = 0.0498 \approx 0.05
\end{equation}

\textbf{4. Complete Calculation:}
\begin{equation}
	2.01 \times 10^{-3} \times 0.0647 = 1.30 \times 10^{-4}
\end{equation}

\textbf{What yields $10^{-4}$:}
It is the T0-calculated Higgs parameter factor $0.0647 \approx 6.5 \times 10^{-2}$ that reduces the loop suppression by factor 20:

\begin{equation}
	2.01 \times 10^{-3} \times 6.5 \times 10^{-2} = 1.3 \times 10^{-4}
\end{equation}

The $10^{-4}$ factor arises from: \textbf{QFT Loop Suppression} ($\sim 10^{-3}$) $\times$ \textbf{T0 Higgs Sector Suppression} ($\sim 10^{-1}$) $=$ $10^{-4}$.

	\section{The Mass Scaling Exponent $\kappa$}
	
	From the fractal dimension follows directly:
	
	\begin{equation}
		\kappa = \frac{D_f^{\,\text{eff}}}{2} = \frac{2.973}{2} \approx 1.487
	\end{equation}
	
	This exponent determines the nonlinear mass scaling in T0-theory.
	
	\section{Lepton Masses from Quantum Numbers}
	
	The masses of leptons follow from the fundamental mass formula:
	
	\begin{equation}
		m_x = \frac{\hbar c}{\xi^2} \times f(n, l, j)
	\end{equation}
	
	where $f(n, l, j)$ is a function of quantum numbers:
	
	\begin{align}
		f(n, l, j) = \sqrt{n(n+l)} \times \left[j + \frac{1}{2}\right]^{1/2}
	\end{align}
	
	For the three leptons we obtain:
	
	\begin{itemize}
		\item Electron $(n=1, l=0, j=1/2)$: $m_e = 0.511$ MeV
		\item Muon $(n=2, l=0, j=1/2)$: $m_\mu = 105.66$ MeV
		\item Tau $(n=3, l=0, j=1/2)$: $m_\tau = 1776.86$ MeV
	\end{itemize}
	
	As claimed in this document, these masses are not empirical inputs but follow from $\xi$ and quantum numbers (for the classification, see the addendum, R77).
	
	\section{The Characteristic Energy $E_0$}
	
	The characteristic energy $E_0$ follows from the geometric mean of the electron and muon masses with fractal correction (following section): $E_0 = \sqrt{m_e m_\mu/K_{\text{frac}}} = 7.398$ MeV (Doc.~190, R72).
	
	\textit{(Corrected on 30~Sept.~2026: The gravitational-geometric derivation via $E_0^2 = \beta_T\,yv/r_g^2$ with $r_g = 2Gm_\mu$ and $G = \xi^2/(4m_\mu)$, i.e.\ $E_0^2 = 4\sqrt{2}\,m_\mu/\xi^4$, has been removed. It gives $1.4 \times 10^{9}$ MeV instead of 7.398 MeV and is dimensionally inconsistent (already noted on 29~Sept.~2026); no derivation supporting it exists in the corpus.)}
	
	\section{Alternative Derivation of $E_0$ from Mass Ratios}
	
	\subsection{The Geometric Mean of Lepton Energies}
	
	An alternative derivation of $E_0$ results directly from the geometric mean of electron and muon masses:
	
	\begin{equation}
		E_0 = \sqrt{m_e \cdot m_\mu} \cdot c^2
	\end{equation}
	
	With the masses calculated from quantum numbers:
	\begin{align}
		E_0 &= \sqrt{0.511 \text{ MeV} \times 105.66 \text{ MeV}}\\
		&= \sqrt{54.00 \text{ MeV}^2}\\
		&= 7.35 \text{ MeV}
	\end{align}
	
	\subsection{Comparison with the Value Used}
	
	The value from the geometric mean (7.35 MeV) agrees well with the value used (7.398 MeV) \textit{(corrected on 29~Sept.~2026: previously ``the value from gravitational derivation'' -- that derivation is not tenable, see above)}. The difference is less than 1\%:
	
	\begin{equation}
		\Delta = \frac{7.398 - 7.35}{7.35} \times 100\% = 0.65\%
	\end{equation}
	
	\subsection{Physical Interpretation}
	
	That $E_0$ corresponds to the geometric mean of the lepton energies is interpreted physically as follows:
	
	\begin{itemize}
		\item $E_0$ represents a natural electromagnetic energy scale between electron and muon
		\item The relationship is purely geometric and requires no knowledge of $\alpha$
		\item The mass ratio $m_\mu/m_e = 206.77$ is itself determined by quantum numbers
	\end{itemize}
	
	\subsection{Precision Correction}
	
	The small difference between 7.35 MeV and 7.398 MeV is the fractal correction (Doc.~190, R72):
	
	\begin{equation}
		E_0^{\text{corrected}} = \frac{E_0^{\text{geom}}}{\sqrt{K_{\text{frac}}}} = \frac{7.348}{\sqrt{1 - 100\xi}} = \frac{7.348}{0.99331} = 7.397 \approx 7.398 \text{ MeV}
	\end{equation}
	
	\textit{(Corrected on 30~Sept.~2026: previously $E_0^{\text{geom}} \times (1 + \alpha/2\pi) = 7.35 \times 1.00116 = 7.358$~MeV and \enquote{additional higher-order quantum corrections are expected to bring the value to 7.398 MeV} -- $\alpha/2\pi$ is not the fractal factor; the gap is closed completely by $1/\sqrt{K_{\text{frac}}} = 1.00673$: $7.348 \times 1.00673 = 7.397$.)}
	
	\subsection{Verification of Fine Structure Constant}
	
	With the geometrically derived $E_0 = 7.35$ MeV:
	
	\begin{align}
		\varepsilon &= \xi \cdot E_0^2\\
		&= (1.333 \times 10^{-4}) \times (7.35)^2\\
		&= (1.333 \times 10^{-4}) \times 54.02\\
		&= 7.20 \times 10^{-3}\\
		&= \frac{1}{138.9}
	\end{align}
	
	With the fractally corrected value $E_0 = 7.348/\sqrt{K_{\text{frac}}}$ one gets $\alpha^{-1} = K_{\text{frac}}/(\xi \cdot 7.348^2) = 0.98667 \times 138.91 = 137.06$ ($+0.017\,\%$). \textit{(Corrected on 30~Sept.~2026: previously \enquote{The remaining deviation from $1/137.036$ is expected to vanish in the more precise calculation with corrected values} -- the correction is $K_{\text{frac}}$, see above.)} This supports the statement that $E_0$ can be derived independently of knowledge of the fine structure constant.
	
	\textit{(Corrected on 30~Sept.~2026: The section \enquote{Two Geometric Paths to $E_0$: Consistency Check} has been removed. It placed $E_0^2 = 4\sqrt{2}\,m_\mu/\xi^4$ as a first path beside the geometric mean, equated the two ($m_e = 4\sqrt{2}/\xi^4 = 1.79\times10^{16}$) and read from this a \enquote{duality} and \enquote{overdetermination} of $E_0$. That path is not tenable (see above). The two routes to $E_0$ are the geometric mean without and with fractal correction, $7.348$ and $7.397$ MeV (previous section); conversely, the measured $\alpha$ fixes the factor $K_{\text{frac}} = 0.98650$, $1.7\times10^{-4}$ from $1 - 100\xi$ (A130).)}
	
	\section{The T0 Coupling Parameter $\varepsilon$}
	
	The T0 coupling parameter results as:
	
	\begin{equation}
		\varepsilon = \xi \cdot E_0^2
	\end{equation}
	
	With the derived values:
	\begin{align}
		\varepsilon &= (1.333 \times 10^{-4}) \times (7.398 \text{ MeV})^2\\
		&= 7.297 \times 10^{-3}\\
		&= \frac{1}{137.036}
	\end{align}
	
	The agreement with the fine structure constant was not presupposed but emerges as a result of the geometric derivation.
\section*{The Simplest Formula for the Fine-Structure Constant}

\[
\boxed{\alpha = \xi \cdot \left(\frac{E_0}{1 \text{ MeV}}\right)^2}
\]
\begin{tcolorbox}[colback=red!5!white,colframe=red!75!black]
	\textbf{Important:} The normalization $(1 \text{ MeV})^2$ is required for dimensionless results.
\end{tcolorbox}	
	\section{Alternative Derivation via Fractal Correction}
	
	As a second path, $\alpha$ can also be derived through fractal renormalization:
	
	\begin{equation}
		\alpha_{\text{bare}}^{-1} = 3\pi \times \xi^{-1} \times \ln\left(\frac{\Lambda_{\text{Planck}}}{m_\mu}\right)
	\end{equation}
	
	With the fractal damping factor:
	\begin{equation}
		D_{\text{frac}} = \left(\frac{\lambda_C^{(\mu)}}{\ell_P}\right)^{D_f-2} = 4.2 \times 10^{-5}
	\end{equation}
	
	we obtain:
	\begin{equation}
		\alpha^{-1} = \alpha_{\text{bare}}^{-1} \times D_{\text{frac}} = 137.036
	\end{equation}
	
	This derivation leads to the same result. \textit{(Note of 30~Sept.~2026: The chain does not yield the stated value. $\lambda_C^{(\mu)}/\ell_P = 1.156 \times 10^{20}$; raised to $D_f - 2$ with $D_f = 3 - \xi$ this gives about $1.15 \times 10^{20}$, not $4.2 \times 10^{-5}$. With $\alpha_{\text{bare}}^{-1} = 3\pi \times 7500 \times \ln(1.2209 \times 10^{22}/105.66) = 3\pi \times 7500 \times 46.196 = 3.265 \times 10^{6}$, the value $4.2 \times 10^{-5}$ is exactly $137.036/3.265 \times 10^{6}$, i.e.\ fitted to $\alpha$ (the exponent would have to be $-0.218$). This is not an independent derivation; the load-bearing relation is $\alpha = \xi E_0^2$ with $E_0 = 7.348/\sqrt{K_{\text{frac}}}$.)}
	
	\section{Clarification: The Two Different $\kappa$ Parameters}
	
	\subsection{Important Distinction}
	
	In T0-theory literature, two physically different parameters are denoted by the symbol $\kappa$, which can lead to confusion. These must be clearly distinguished:
	
	\begin{enumerate}
		\item $\kappa_{\text{mass}} = 1.487$ - The fractal mass scaling exponent
		\item $\kappa_{\text{grav}}$ - The gravitational field parameter
	\end{enumerate}
	
	\subsection{The Mass Scaling Exponent $\kappa_{\text{mass}}$}
	
	This parameter was already derived in Section 4:
	
	\begin{equation}
		\kappa_{\text{mass}} = \frac{D_f^{\,\text{eff}}}{2} = 1.487
	\end{equation}
	\textit{(Corrected on 30~Sept.~2026: previously $\frac{D_f}{2} = 1.47$ -- by Section~4, $D_f^{\text{eff}}/2 = 2.973/2 = 1.4865$; $1.47$ follows from no $D_f$ value ($D_f = 3-\xi$ would give $1.49993$). Likewise in the list above and in the symbol table.)}
	
	It is dimensionless and determines the scaling in the formula for magnetic moments:
	
	\begin{equation}
		a_x \propto \left(\frac{m_x}{m_\mu}\right)^{\kappa_{\text{mass}}}
	\end{equation}
	
	\subsection{The Gravitational Field Parameter $\kappa_{\text{grav}}$}
	
	This parameter arises from the coupling between the intrinsic time field and matter. The T0 Lagrangian density reads:
	
	\begin{equation}
		\mathcal{L}_{\text{intrinsic}} = \frac{1}{2}\partial_\mu T \partial^\mu T - \frac{1}{2}T^2 - \frac{\rho}{T}
	\end{equation}
	
	The resulting field equation:
	
	\begin{equation}
		\nabla^2 T = -\frac{\rho}{T^2}
	\end{equation}
	
	leads to a modified gravitational potential:
	
	\begin{equation}
		\Phi(r) = -\frac{GM}{r} + \kappa_{\text{grav}} r
	\end{equation}
	
	\subsection{Relationship Between $\kappa_{\text{grav}}$ and Fundamental Parameters}
	
	In natural units:
	
	\begin{equation}
		\kappa_{\text{grav}}^{\text{nat}} = \beta_T^{\text{nat}} \cdot \frac{yv}{r_g^2}
	\end{equation}
	
	With $\beta_T = 1$ and $r_g = 2Gm_\mu$:
	
	\begin{equation}
		\kappa_{\text{grav}} = \frac{y_\mu \cdot v}{(2Gm_\mu)^2} = \frac{\sqrt{2} m_\mu \cdot v}{v \cdot 4G^2m_\mu^2} = \frac{\sqrt{2}}{4G^2m_\mu}
	\end{equation}
	
	\subsection{Numerical Value and Physical Significance}
	
	In SI units:
	
	\begin{equation}
		\kappa_{\text{grav}}^{\text{SI}} \approx 4.8 \times 10^{-11} \text{ m/s}^2
	\end{equation}
	
	This linear term in the gravitational potential:
	\begin{itemize}
		\item Is intended to explain observed flat rotation curves of galaxies \textbf{[S]}
		\item Would thereby make dark matter unnecessary \textbf{[S]}
		\item Arises naturally from time field-matter coupling
	\end{itemize}
	
\section{Complete Mapping: Standard Model Parameters to T0 Correspondences}
\label{sec:sm_t0_mapping}

\subsection{Overview of Parameter Reduction}
\label{subsec:parameter_overview}

The Standard Model requires over 20 free parameters that must be determined experimentally. The T0 system traces these back to the one parameter $\xi$, supplemented by integer structural choices (rational prefactors, exponents) that enter as motivated settings, and by reference anchors such as $v$ and $E_P$ for unit conversion (for the classification of the early formulas, see the addendum, R77):

\begin{equation}
	\boxed{\xi = \frac{4}{3} \times 10^{-4}}
\end{equation}

\subsection{Hierarchically Ordered Parameter Mapping Table}
\label{subsec:hierarchical_mapping}

The table is organized so that each parameter is defined before being used in subsequent formulas.

\begin{table}[htbp]
\centering
\caption{Standard Model Parameters in T0 Hierarchy — Part 1: Levels 0--3}
\label{tab:sm_hierarchy_part1}
\footnotesize
\adjustbox{max width=\textwidth}{%
\begin{tabular}{p{3.8cm}p{2.5cm}p{3.5cm}p{2.8cm}}
\toprule
\textbf{SM Parameter} & \textbf{SM Value} & \textbf{T0 Formula} & \textbf{T0 Value} \\
\midrule
	% LEVEL 0: FUNDAMENTAL CONSTANT
	\multicolumn{4}{l}{\textbf{LEVEL 0: FUNDAMENTAL GEOMETRIC CONSTANT}} \\
	\midrule
	
	Geometric parameter $\xi$ & -- & $\xi = \frac{4}{3} \times 10^{-4}$ & $1.333 \times 10^{-4}$ \\
	& & (from geometric) & (exact) \\[0.3em]
	
	\midrule
	% LEVEL 1: DIRECT DERIVATIVES FROM XI
	\multicolumn{4}{l}{\textbf{LEVEL 1: PRIMARY COUPLING CONSTANTS (dependent only on $\xi$)}} \\
	\midrule
	
	Strong coupling $\alpha_S$ & $\alpha_S \approx 0.118$ & $\alpha_S = \xi^{-1/3}$ & $19.57$ \textit{(Corrected on 30~Sept.~2026: previously $9.65$ -- $(1.333 \times 10^{-4})^{-1/3} = 19.57$.)} \\
	& (at $M_Z$) & $= (1.333 \times 10^{-4})^{-1/3}$ & (nat. units) \\[0.3em]
	
	Weak coupling $\alpha_W$ & $\alpha_W \approx 1/30$ & $\alpha_W = \xi^{1/2}$ & $1.15 \times 10^{-2}$ \\
	& & $= (1.333 \times 10^{-4})^{1/2}$ & \\[0.3em]
	
	Gravitational coupling $\alpha_G$ & not in SM & $\alpha_G = \xi^{2}$ & $1.78 \times 10^{-8}$ \\
	& & $= (1.333 \times 10^{-4})^{2}$ & \\[0.3em]
	
	Electromagnetic coupling & $\alpha = 1/137.036$ & $\alpha_{EM} = 1$ (convention) & $1$ \\
	& & $\varepsilon_T = \xi \cdot E_0^2$ & $7.297 \times 10^{-3}$ \textit{(Corrected on 30~Sept.~2026: previously $\varepsilon_T = \xi \cdot \sqrt{3/(4\pi^2)} = 3.7 \times 10^{-5}$ -- contradicted $\varepsilon_T = \xi \cdot E_0^2$ in the text and in the symbol list.)} \\
	& & (physical coupling) & (*see note) \\[0.3em]
	
	\midrule
	% LEVEL 2: ENERGY SCALES
	\multicolumn{4}{l}{\textbf{LEVEL 2: ENERGY SCALES (dependent on $\xi$ and Planck scale)}} \\
	\midrule
	
	Planck energy $E_P$ & $1.22 \times 10^{19}$ GeV & Reference scale & $1.22 \times 10^{19}$ GeV \\
	& & (from $G, \hbar, c$) & \\[0.3em]
	
Higgs-VEV $v$ & $246.22$ GeV & $v = \xi^4 E_P/(5\pi)$ & $245.65$ GeV \textit{(Corrected on 30~Sept.~2026: previously $v = \frac{4}{3}\xi_0^{-1/2} K_{\text{quantum}}$ with $246.2$ -- $K_{\text{quantum}}$ is fitted to $v$ and the appendix chain is faulty; structural formula Docs.~149, 383.)} \\
& (theoretisch) & (see appendix) & \\[0.3em]
	
	QCD scale $\Lambda_{QCD}$ & $\sim 217$ MeV & $\Lambda_{QCD} = v \cdot \xi^{1/3}$ & $12.6$ GeV \\
	& (free parameter) & $= 246 \text{ GeV} \cdot \xi^{1/3}$ & \\[0.3em]
	
	\midrule
	% LEVEL 3: HIGGS SECTOR
	\multicolumn{4}{l}{\textbf{LEVEL 3: HIGGS SECTOR (dependent on $v$)}} \\
	\midrule
	
	Higgs mass $m_h$ & $125.25$ GeV & $m_h = v \cdot \xi^{1/4}$ & $26.4$ GeV \\
	& (measured) & $= 246 \cdot (1.333 \times 10^{-4})^{1/4}$ & \\[0.3em]
	
	Higgs self-coupling $\lambda_h$ & $0.13$ & $\lambda_h = \frac{m_h^2}{2v^2}$ & $0.129$ \\
	& (derived) & $= \frac{(125)^2}{2(246)^2}$ & \\[0.3em]
	
\end{tabular}%
}
\end{table}
\textit{(Corrected on 30~Sept.~2026: $\Lambda_{QCD} = v\cdot\xi^{1/3}$ previously $200$~MeV -- $246\times0.05109 = 12.57$~GeV; $m_h = v\cdot\xi^{1/4}$ previously $125$~GeV -- $246\times0.10746 = 26.4$~GeV. Neither formula reproduces the measured values ($\sim 217$~MeV and $125.25$~GeV); $\lambda_h$ below it is computed with the measured $m_h = 125$~GeV.)}

\begin{table}[htbp]
\centering
\caption{Standard Model Parameters in T0 Hierarchy — Part 2: Levels 4--5}
\label{tab:sm_hierarchy_part2}
\footnotesize
\adjustbox{max width=\textwidth}{%
\begin{tabular}{p{3.8cm}p{2.5cm}p{3.5cm}p{2.8cm}}
\toprule
\textbf{SM Parameter} & \textbf{SM Value} & \textbf{T0 Formula} & \textbf{T0 Value} \\
\midrule
	\midrule
	% LEVEL 4: FERMION MASSES
	\multicolumn{4}{l}{\textbf{LEVEL 4: FERMION MASSES (dependent on $v$ and $\xi$)}} \\
	\midrule
	
	\multicolumn{4}{l}{\textit{Leptons:}} \\
	
	Electron mass $m_e$ & $0.511$ MeV & $m_e = v \cdot \frac{4}{3} \cdot \xi^{3/2}$ & $0.502$ MeV \\
	& (free parameter) & $= 246 \text{ GeV} \cdot \frac{4}{3} \cdot \xi^{3/2}$ & \\[0.3em]
	
	Muon mass $m_\mu$ & $105.66$ MeV & $m_\mu = v \cdot \frac{16}{5} \cdot \xi^1$ & $105.0$ MeV \\
	& (free parameter) & $= 246 \text{ GeV} \cdot \frac{16}{5} \cdot \xi$ & \\[0.3em]
	
	Tau mass $m_\tau$ & $1776.86$ MeV & $m_\tau = v \cdot \frac{25}{9} \cdot \xi^{2/3}$ & $1783$ MeV \\
	& (free parameter) & $= 246 \text{ GeV} \cdot \frac{25}{9} \cdot \xi^{2/3}$ & \\[0.3em]
	
	\multicolumn{4}{l}{\textit{Up-type quarks:}} \\
	
	Up quark mass $m_u$ & $2.16$ MeV & $m_u = v \cdot 6 \cdot \xi^{3/2}$ & $2.27$ MeV \\
	
	Charm quark mass $m_c$ & $1.27$ GeV & $m_c = v \cdot 2 \cdot \xi^{2/3}$ & $1.284$ GeV \\
	
	Top quark mass $m_t$ & $172.76$ GeV & $m_t = v \cdot \frac{1}{28} \cdot \xi^{-1/3}$ & $173.0$ GeV \\
	
	\multicolumn{4}{l}{\textit{Down-type quarks:}} \\
	
	Down quark mass $m_d$ & $4.67$ MeV & $m_d = v \cdot \frac{25}{2} \cdot \xi^{3/2}$ & $4.72$ MeV \\
	
	Strange quark mass $m_s$ & $93.4$ MeV & $m_s = v \cdot 3 \cdot \xi^1$ & $98.4$ MeV \textit{(Corrected on 1~Oct.~2026: previously $97.9$~MeV -- $246~\text{GeV} \cdot 3 \cdot \xi = 98.4$~MeV ($+5.4\,\%$ against $93.4$~MeV); Doc.~006 uses $r_s = 26/9$ and hence $94.76$~MeV.)} \\
	
	Bottom quark mass $m_b$ & $4.18$ GeV & $m_b = v \cdot \frac{3}{2} \cdot \xi^{1/2}$ & $4.254$ GeV \\
	
	\midrule
	% LEVEL 5: NEUTRINO MASSES
	\multicolumn{4}{l}{\textbf{LEVEL 5: NEUTRINO MASSES (dependent on $v$ and double $\xi$)}} \\
	\midrule
	
	Electron neutrino $m_{\nu_e}$ & $< 2$ eV & $m_{\nu_e} = v \cdot r_{\nu_e} \cdot \xi^{3/2} \cdot \xi^3$ & $\sim 10^{-6}$ eV \\
	& (upper limit) & with $r_{\nu_e} \sim 1$ & (prediction) \\[0.3em]
	
	Muon neutrino $m_{\nu_\mu}$ & $< 0.19$ MeV & $m_{\nu_\mu} = v \cdot r_{\nu_\mu} \cdot \xi^{1} \cdot \xi^3$ & $\sim 10^{-4}$ eV \\
	
	Tau neutrino $m_{\nu_\tau}$ & $< 18.2$ MeV & $m_{\nu_\tau} = v \cdot r_{\nu_\tau} \cdot \xi^{2/3} \cdot \xi^3$ & $\sim 10^{-3}$ eV \\
	
\end{tabular}%
}
\end{table}
\textit{(Corrected on 30~Sept.~2026: tau previously $m_\tau = v\cdot\frac54\cdot\xi^{2/3} = 1778$~MeV -- with $\frac54$ one gets $802.5$~MeV, $1783$~MeV only with $r_\tau = \frac{25}{9}$ (Docs.~006, 046); charm previously $m_c = v\cdot\frac89\cdot\xi^{2/3} = 1.279$~GeV -- with $\frac89$ one gets $0.571$~GeV, $1.284$~GeV only with $r_c = 2$ (Doc.~046); neutrinos previously $\sim10^{-3}$, $\sim10^{-2}$, $\sim10^{-1}$~eV -- with $r\sim1$: $v\,\xi^{3/2}\xi^3 = 9.0\times10^{-7}$~eV, $v\,\xi\,\xi^3 = 7.8\times10^{-5}$~eV, $v\,\xi^{2/3}\xi^3 = 1.5\times10^{-3}$~eV.)}

\begin{table}[htbp]
\centering
\caption{Standard Model Parameters in T0 Hierarchy — Part 3: Levels 6--7}
\label{tab:sm_hierarchy_part3}
\footnotesize
\adjustbox{max width=\textwidth}{%
\begin{tabular}{p{3.8cm}p{2.5cm}p{3.5cm}p{2.8cm}}
\toprule
\textbf{SM Parameter} & \textbf{SM Value} & \textbf{T0 Formula} & \textbf{T0 Value} \\
\midrule
	\midrule
	% LEVEL 6: MIXING PARAMETERS
	\multicolumn{4}{l}{\textbf{LEVEL 6: MIXING MATRICES (dependent on mass ratios)}} \\
	\midrule
	
	\multicolumn{4}{l}{\textit{CKM Matrix (Quarks):}} \\
	
	$|V_{us}|$ (Cabibbo) & $0.22452$ & $|V_{us}| = \sqrt{\frac{m_d}{m_s}} \cdot f_{Cab}$ & $0.21$ \textit{(Corrected on 1~Oct.~2026: previously $0.225$ -- with the T0 masses of this table ($m_d = 4.72$, $m_s = 97.9$~MeV) the formula gives $0.209$, with Doc.~006 ($4.734$/$94.76$~MeV) $0.213$; it does not reproduce the SM value ($-5\,\%$).)} \\
	& & with $f_{Cab} = \sqrt{\frac{m_s - m_d}{m_s + m_d}}$ & \\[0.3em]
	
	$|V_{ub}|$ & $0.00365$ & $|V_{ub}| = \sqrt{\frac{m_d}{m_b}} \cdot \xi^{1/4}$ & $0.0037$ \\
	
	$|V_{ud}|$ & $0.97446$ & $|V_{ud}| = \sqrt{1 - |V_{us}|^2 - |V_{ub}|^2}$ & $0.974$ \\
	& & (unitarity) & \\[0.3em]
	
	CKM CP phase $\delta_{CKM}$ & $1.20$ rad & $\delta_{CKM} = \arcsin(2\sqrt{2}\xi^{1/2}/3)$ & $0.0109$ rad \textit{(Corrected on 1~Oct.~2026: previously $1.2$ rad -- $\arcsin(2\sqrt2\cdot\sqrt{\xi}/3) = \arcsin(0.0109) = 0.0109$ rad; only without the factor $\sqrt{\xi}$ would one get $\arcsin(2\sqrt2/3) = 1.231$ rad.)} \\
	
	\multicolumn{4}{l}{\textit{PMNS Matrix (Neutrinos):}} \\
	
	$\theta_{12}$ (Solar) & $33.44°$ & $\theta_{12} = \arcsin\sqrt{m_{\nu_1}/m_{\nu_2}}$ & $33.5°$ \\
	
	$\theta_{23}$ (Atmospheric) & $49.2°$ & $\theta_{23} = \arcsin\sqrt{m_{\nu_2}/m_{\nu_3}}$ & $49°$ \\
	
	$\theta_{13}$ (Reactor) & $8.57°$ & $\theta_{13} = \arcsin(\xi^{1/3})$ & $2.93°$ \textit{(Corrected on 1~Oct.~2026: previously $8.6°$ -- $\xi^{1/3} = 0.0511$, $\arcsin(0.0511) = 2.93°$; $8.59°$ only follows with the argument $(25\xi)^{1/3}$, $25\xi = 1/300 = 1/(5|A_5|)$, the corpus formula of Doc.~328/340.)} \\
	
	PMNS CP phase $\delta_{CP}$ & unknown & $\delta_{CP} = \pi(1 - 2\xi)$ & $3.14$ rad \textit{(Corrected on 30~Sept.~2026: previously $1.57$ rad -- $\pi(1 - 2\xi) = 3.1408$ rad.)} \\
	
	\midrule
	% LEVEL 7: DERIVED PARAMETERS
	\multicolumn{4}{l}{\textbf{LEVEL 7: DERIVED PARAMETERS}} \\
	\midrule
	
	Weinberg $\sin^2\theta_W$ & $0.2312$ & $\tfrac{1{-}\sqrt{1{-}4\alpha_W}}{4}$ & $0.0058$ \textit{(Corrected on 30~Sept.~2026: previously $0.231$ -- with $\alpha_W = 0.01155$: $(1 - \sqrt{0.9538})/4 = 0.0058$; the formula does not reproduce the SM value.)} \\
	& & with $\alpha_W$ from Level 1 & \\[0.3em]
	
	Strong CP phase $\theta_{QCD}$ & $< 10^{-10}$ & $\theta_{QCD} = \xi^{2}$ & $1.78 \times 10^{-8}$ \\
	& (upper limit) & & (prediction) \\
	
\end{tabular}%
}
\end{table}
	
	\textit{(Note of 30~Sept.~2026: The prediction $\theta_{QCD} = \xi^2 = 1.78 \times 10^{-8}$ lies about 178 times above the bound $10^{-10}$ from the electric dipole moment of the neutron quoted in the table and is therefore experimentally excluded.)}
\normalsize

\subsection{The Hierarchical Derivation Structure}
\label{subsec:hierarchical_structure}

The table shows the clear hierarchy of parameter derivation:

\begin{enumerate}
	\item \textbf{Level 0}: Only $\xi$ as fundamental constant
	\item \textbf{Level 1}: Coupling constants directly from $\xi$
	\item \textbf{Level 2}: Energy scales from $\xi$ and reference scales
	\item \textbf{Level 3}: Higgs parameters from energy scales
	\item \textbf{Level 4}: Fermion masses from $v$ and $\xi$
	\item \textbf{Level 5}: Neutrino masses with additional suppression
	\item \textbf{Level 6}: Mixing parameters from mass ratios
	\item \textbf{Level 7}: Further derived parameters
\end{enumerate}

Each level uses only parameters that were defined in previous levels.

\subsection{Critical Notes}
\label{subsec:critical_notes}

\textbf{(*) Note on the Fine Structure Constant:}

The fine structure constant has a dual function in the T0 system:
\begin{itemize}
	\item $\alpha_{EM} = 1$ is a \textbf{unit convention} (like $c = 1$)
	\item $\varepsilon_T = \xi \cdot E_0^2 = 7.297 \times 10^{-3}$ is the \textbf{physical EM coupling} \textit{(Corrected on 30~Sept.~2026: previously $\varepsilon_T = \xi \cdot f_{geom}$ -- with $f_{geom} = \sqrt{3/(4\pi^2)}$ this would give $3.7 \times 10^{-5}$, not $\alpha$; the text computes $\varepsilon = \xi E_0^2$.)}
\end{itemize}

\textbf{Unit System:}
All T0 values apply in natural units with $\hbar = c = 1$. Transformation to SI units is required for experimental comparisons.
\section{Cosmological Parameters: Standard Cosmology ($\Lambda$CDM) vs T0 System}
\label{sec:cosmic_t0_mapping}

\subsection{Basic Difference of the Approaches}
\label{subsec:paradigm_shift}

\begin{tcolorbox}[colback=red!5!white,colframe=red!75!black,title=Note: Basic Differences]
	The T0 system postulates a \textbf{static, eternal universe} without a Big Bang, while standard cosmology is based on an \textbf{expanding universe} with a Big Bang. The parameters are therefore often not directly comparable but represent different physical concepts.

	The present corpus deliberately sets the cosmic sector aside, because the Standard Model does not derive it either (P39). The cosmological T0 statements in this section are therefore speculative \textbf{[S]}.
\end{tcolorbox}

\subsection{Hierarchically Ordered Cosmological Parameters}
\label{subsec:cosmic_hierarchical_mapping}

\begin{table}[htbp]
\centering
\caption{Cosmological Parameters in Hierarchical Order — Part 1: Geometric \& Energy Scales}
\label{subsec:cosmic_hierarchical_part1}
\footnotesize
\adjustbox{max width=\textwidth}{%
\begin{tabular}{p{3.8cm}p{2.5cm}p{3.5cm}p{2.8cm}}
\toprule
\textbf{Parameter} & \textbf{$\Lambda$CDM Value} & \textbf{T0 Formula} & \textbf{T0 Interpretation} \\
\midrule
	% LEVEL 0: FUNDAMENTAL CONSTANT
	\multicolumn{4}{l}{\textbf{LEVEL 0: FUNDAMENTAL GEOMETRIC CONSTANT}} \\
	\midrule
	
	Geometric parameter $\xi$ & not present & $\xi = \frac{4}{3} \times 10^{-4}$ & $1.333 \times 10^{-4}$ \\
	& & (from geometric) & basis of all derivations \\[0.3em]
	
	\midrule
	% LEVEL 1: PRIMARY COSMIC PARAMETERS
	\multicolumn{4}{l}{\textbf{LEVEL 1: PRIMARY ENERGY SCALES (dependent only on $\xi$)}} \\
	\midrule
	
	Characteristic energy & -- & $E_\xi = \frac{1}{\xi} = \frac{3}{4} \times 10^{4}$ & $7500$ (nat. units) \\
	& & & CMB energy scale \\[0.3em]
	
	Characteristic length & -- & $L_\xi = \xi$ & $1.33 \times 10^{-4}$ \\
	& & & (nat. units) \\[0.3em]
	
	$\xi$-field energy density & -- & $\rho_\xi = E_\xi^4$ & $3.16 \times 10^{16}$ \\
	& & & vacuum energy density \\[0.3em]
	
	\midrule
	% LEVEL 2: CMB PARAMETERS
	\multicolumn{4}{l}{\textbf{LEVEL 2: CMB PARAMETERS (dependent on $\xi$ and $E_\xi$)}} \\
	\midrule
	
	CMB temperature today & $T_0 = 2.7255$ K & $T_{CMB} = \frac{16}{9} \xi^2 \cdot E_\xi$ & $2.725$ K \\
	& (measured) & $= \frac{16}{9} \cdot (1.33 \times 10^{-4})^2 \cdot 7500$ & (calculated) \\[0.3em]
	
	CMB energy density & $\rho_{CMB} = 4.64 \times 10^{-31}$ kg/m³ & $\rho_{CMB} = \frac{\pi^2}{15} T_{CMB}^4$ & $4.2 \times 10^{-14}$ J/m³ \\
	& & Stefan-Boltzmann & (nat. units) \\[0.3em]
	
	CMB anisotropy & $\Delta T/T \sim 10^{-5}$ & $\delta T = \xi^{1/2} \cdot T_{CMB}$ & $\sim 10^{-5}$ \\
	& (Planck satellite) & quantum fluctuation & (predicted) \\[0.3em]
	
\end{tabular}%
}
\end{table}

\begin{table}[htbp]
\centering
\caption{Cosmological Parameters — Part 2: Redshift \& Dark Components}
\label{subsec:cosmic_hierarchical_part2}
\footnotesize
\adjustbox{max width=\textwidth}{%
\begin{tabular}{p{3.8cm}p{2.5cm}p{3.5cm}p{2.8cm}}
\toprule
\textbf{Parameter} & \textbf{$\Lambda$CDM Value} & \textbf{T0 Formula} & \textbf{T0 Interpretation} \\
\midrule
	\midrule
	% LEVEL 3: REDSHIFT
	\multicolumn{4}{l}{\textbf{LEVEL 3: REDSHIFT (dependent on $\xi$ and wavelength)}} \\
	\midrule
	
	& (Planck 2020) & Static universe & \\[0.3em]
	
	Redshift $z$ & $z = \frac{\Delta\lambda}{\lambda}$ & $z(\lambda, d) = \xi \cdot \lambda \cdot d$ & Fractal path-lengthening \\
	& (expansion) & Wavelength-dependent (original state; cf.\ K6) & not expansion \\[0.3em]
	
	Effective $H_0$ & $66.2$ km/s/Mpc & see Document~026 \\
	(interpreted) & & at $\lambda_{ref} = 550$ nm & (apparent) \\[0.3em]
	
	\midrule
	% LEVEL 4: DARK COMPONENTS
	\multicolumn{4}{l}{\textbf{LEVEL 4: DARK COMPONENTS}} \\
	\midrule
	
	Dark energy $\Omega_\Lambda$ & $0.6847 \pm 0.0073$ & Not required & $0$ \\
	& (68.47\% of universe) & Static universe & eliminated \\[0.3em]
	
	Dark matter $\Omega_{DM}$ & $0.2607 \pm 0.0067$ & $\xi$-field effects & $0$ \\
	& (26.07\% of universe) & Modified gravity & eliminated \\[0.3em]
	
	Baryonic matter $\Omega_b$ & $0.0492 \pm 0.0003$ & All matter & $1.0$ \\
	& (4.92\% of universe) & & (100\%) \\[0.3em]
	
	Cosmological constant $\Lambda$ & $(1.1 \pm 0.02) \times 10^{-52}$ m$^{-2}$ & $\Lambda = 0$ & $0$ \\
	& & No expansion & eliminated \\[0.3em]
	
\end{tabular}%
}
\end{table}

\begin{table}[htbp]
\centering
\caption{Cosmological Parameters — Part 3: Universe Structure, Predictions \& Energy Balances}
\label{subsec:cosmic_hierarchical_part3}
\footnotesize
\adjustbox{max width=\textwidth}{%
\begin{tabular}{p{3.8cm}p{2.5cm}p{3.5cm}p{2.8cm}}
\toprule
\textbf{Parameter} & \textbf{$\Lambda$CDM Value} & \textbf{T0 Formula} & \textbf{T0 Interpretation} \\
\midrule
	\midrule
	% LEVEL 5: UNIVERSE AGE AND STRUCTURE
	\multicolumn{4}{l}{\textbf{LEVEL 5: UNIVERSE STRUCTURE}} \\
	\midrule
	
	Universe age & $13.787 \pm 0.020$ Gyr & $t_{univ} = \infty$ & Eternal \\
	& (since Big Bang) & No beginning/end & Static \\[0.3em]
	
	Big Bang & $t = 0$ & No Big Bang & -- \\
	& Singularity & Excluded by $T\cdot E = 1$ (Doc.~306) & Excluded \\[0.3em]
	
	Decoupling (CMB) & $z \approx 1100$ ($\Lambda$CDM) & CMB from $\xi$-field & Continuous \\
	& $t = 380,000$ years & Vacuum fluctuation & generation \\[0.3em]
	
	Structure formation & Bottom-up & Continuous & Cyclic \\
	& (small → large) & $\xi$-driven & regenerating \\[0.3em]
	
	\midrule
	% LEVEL 6: PREDICTIONS AND TESTS
	\multicolumn{4}{l}{\textbf{LEVEL 6: DISTINGUISHABLE PREDICTIONS}} \\
	\midrule
	
	Hubble tension & Unsolved & Interpreted via & No tension [S] \\
	& $H_0^{local} \neq H_0^{CMB}$ & $\xi$-effects & $H_0^{eff} = 67.45$ \\[0.3em]
	
	JWST early galaxies & Problem & No problem & Expected in \\
	& (formed too early) & Eternal universe & static universe \\[0.3em]
	
	$\lambda$-dependent $z$ & $z$ independent of $\lambda$ & $z \propto \lambda$ & At the limit \\
	& All $\lambda$ same $z$ & $z_{UV} > z_{radio}$ & of testability* \\[0.3em]
	
	Casimir effect & Quantum fluctuation & $F_{Cas} = -\frac{\pi^2}{240} \frac{\hbar c}{d^4}$ & $\xi$-field \\
	& & from $\xi$-geometry & manifestation \\[0.3em]
	
	\midrule
	% LEVEL 7: ENERGY CONSERVATION
	\multicolumn{4}{l}{\textbf{LEVEL 7: ENERGY BALANCES}} \\
	\midrule
	
	Total energy & Not conserved & $E_{total} = const$ & Strictly conserved \\
	& (expansion) & & \\[0.3em]
	
	Mass-energy & $E = mc^2$ & $E = mc^2$ & Identical** \\
	equivalence & & & (see note) \\[0.3em]
	
	Vacuum energy & Problem & $\rho_{vac} = \rho_\xi$ & Naturally from \\
	& ($10^{120}$ discrepancy) & Calculable from $\xi$ & $\xi$ \\[0.3em]
	
	Entropy & Grows monotonically & $S_{total} = const$ & Cyclically \\
	& (heat death) & Regeneration & conserved \\[0.3em]
	
\end{tabular}%
}
\end{table}

\normalsize

\subsection{Critical Differences and Test Possibilities}
\label{subsec:critical_differences}

\begin{table}[h]
	\centering
	\adjustbox{max width=\textwidth}{%
\begin{tabular}{p{4cm}p{5cm}p{5cm}}
		\toprule
		\textbf{Phenomenon} & \textbf{$\Lambda$CDM Explanation} & \textbf{T0 Interpretation [S]} \\
		\midrule
		Redshift & Space expansion & Fractal path-lengthening in the $\xi$-field (geometric path stretching; no photon energy loss) \\
		CMB & Recombination at $z=1100$ ($\Lambda$CDM value) & $\xi$-field equilibrium radiation \\
		Dark energy & 68\% of universe & Not required \\
		Dark matter & 26\% of universe & $\xi$-field gravity effects \\
		Hubble tension & Unsolved (4.4$\sigma$) & Interpreted via $\xi$-effects \\
		JWST paradox & Unexplained early galaxies & No problem in eternal universe \\
		\bottomrule
	\end{tabular}%
}
	\caption{Differences between $\Lambda$CDM and T0 (T0 column speculative, cf.\ P39)}
\end{table}

\subsection{Philosophical Implications}
\label{subsec:philosophical_implications}

In the T0 picture, the following would follow \textbf{[S]}:
\begin{enumerate}
	\item \textbf{Eternal universe}: No beginning, no end -- the question of a first beginning does not arise
	\item \textbf{No singularities}: A minimal length scale $L_0$ (set by $\xi$) replaces the formal Big Bang singularity by a tiny but finite core; the justification via Heisenberg uncertainty has been replaced by $T\cdot E = 1$ (Doc.~306)
	\item \textbf{Energy conservation}: Strictly preserved, no violation through expansion
	\item \textbf{Simplicity}: One constant instead of 6+ parameters (the cosmological sector still requires its own scale, P39)
	\item \textbf{Testability}: Clear, measurable predictions
\end{enumerate}
\section{Appendix: Purely Theoretical Derivation of Higgs VEV from Quantum Numbers}

\subsection{Fundamental theoretical foundations}

\subsubsection{Quantum numbers of leptons in T0 theory}

T0 theory assigns quantum numbers $(n, l, j)$ to each particle, arising from the solution of the three-dimensional wave equation in the energy field:

\textbf{Electron (1st generation):}
\begin{itemize}
	\item Principal quantum number: $n = 1$
	\item Orbital angular momentum: $l = 0$ (s-like, spherically symmetric)
	\item Total angular momentum: $j = 1/2$ (fermion)
\end{itemize}

\textbf{Muon (2nd generation):}
\begin{itemize}
	\item Principal quantum number: $n = 2$
	\item Orbital angular momentum: $l = 1$ (p-like, dipole structure)
	\item Total angular momentum: $j = 1/2$ (fermion)
\end{itemize}

\subsubsection{Universal mass formulas}

T0 theory provides two equivalent formulations for particle masses:

\textbf{Direct method:}
\begin{equation}
	m_i = \frac{1}{\xi_i} = \frac{1}{\xi_0 \times f(n_i, l_i, j_i)}
	\label{eq:direct_mass_formula}
\end{equation}

\textbf{Extended Yukawa method:}
\begin{equation}
	m_i = y_i \times v
	\label{eq:yukawa_mass_formula}
\end{equation}

where:
\begin{itemize}
	\item $\xi_0 = \frac{4}{3} \times 10^{-4}$: Universal geometric parameter
	\item $f(n_i, l_i, j_i)$: Geometric factors from quantum numbers
	\item $y_i$: Yukawa couplings
	\item $v$: Higgs VEV (target quantity)
\end{itemize}

\subsection{Theoretical calculation of geometric factors}

\subsubsection{Geometric factors from quantum numbers}

The geometric factors result from the analytical solution of the three-dimensional wave equation. For the fundamental leptons:

\textbf{Electron $(n=1, l=0, j=1/2)$:}

The ground state solution of the 3D wave equation yields the simplest geometric factor:
\begin{equation}
	f_e(1,0,1/2) = 1
\end{equation}

This is the reference configuration (ground state).

\textbf{Muon $(n=2, l=1, j=1/2)$:}

For the first excited configuration with dipole character, the solution yields:
\begin{equation}
	f_\mu(2,1,1/2) = \frac{16}{5}
\end{equation}

This factor accounts for:
\begin{itemize}
	\item $n^2 = 4$ (energy level scaling)
	\item $\frac{4}{5}$ ($l=1$ dipole correction vs. $l=0$ spherical)
\end{itemize}

\subsubsection{Verification of factors}

The geometric factors must be consistent with the universal T0 structure:

\begin{align}
	\xi_e &= \xi_0 \times f_e = \frac{4}{3} \times 10^{-4} \times 1 = \frac{4}{3} \times 10^{-4}\\
	\xi_\mu &= \xi_0 \times f_\mu = \frac{4}{3} \times 10^{-4} \times \frac{16}{5} = \frac{64}{15} \times 10^{-4}
\end{align}

\subsection{Derivation of mass ratios}

\subsubsection{Theoretical electron-muon mass ratio}

With the geometric factors, it follows from the direct method:

\begin{align}
	\frac{m_\mu}{m_e} &= \frac{\xi_e}{\xi_\mu} = \frac{f_e}{f_\mu} = \frac{1}{\frac{16}{5}} = \frac{5}{16}
\end{align}

\textbf{Note:} This is the inverse ratio. Since $\xi \propto 1/m$, we obtain:

\begin{align}
	\frac{m_\mu}{m_e} &= \frac{f_\mu}{f_e} = \frac{\frac{16}{5}}{1} = \frac{16}{5} = 3.2
\end{align}
\textit{(Note of 30~Sept.~2026: the direct method gives $m_\mu/m_e = 5/16$, or $16/5 = 3.2$ after inversion, the Yukawa method below gives $207.8$; the formulations called \enquote{equivalent} above thus yield different ratios and are not equivalent. In addition, the muon is assigned $l = 1$ here but $l = 0$ in the section \enquote{Lepton Masses from Quantum Numbers}.)}

\subsubsection{Correction through Yukawa couplings}

The Yukawa method accounts for additional quantum field theoretical corrections:

\textbf{Electron:}
\begin{equation}
	y_e = \frac{4}{3} \times \xi^{3/2} = \frac{4}{3} \times \left(\frac{4}{3} \times 10^{-4}\right)^{3/2}
\end{equation}

\textbf{Muon:}
\begin{equation}
	y_\mu = \frac{16}{5} \times \xi^1 = \frac{16}{5} \times \frac{4}{3} \times 10^{-4}
\end{equation}

\subsubsection{Calculation of corrected ratio}

\begin{align}
	\frac{y_\mu}{y_e} &= \frac{\frac{16}{5} \times \frac{4}{3} \times 10^{-4}}{\frac{4}{3} \times \left(\frac{4}{3} \times 10^{-4}\right)^{3/2}}\\
	&= \frac{\frac{16}{5} \times \frac{4}{3} \times 10^{-4}}{\frac{4}{3} \times \frac{4}{3} \times 10^{-4} \times \sqrt{\frac{4}{3} \times 10^{-4}}}\\
	&= \frac{\frac{16}{5}}{\frac{4}{3} \times \sqrt{\frac{4}{3} \times 10^{-4}}}\\
	&= \frac{\frac{16}{5}}{\frac{4}{3} \times 0.01155}\\
	&= \frac{3.2}{0.0154} = 207.8
\end{align}

This theoretical ratio of $207.8$ is very close to the experimental value of $206.768$.

\subsection{Higgs VEV}

In the corpus the Higgs VEV follows from the structural formula
\begin{equation}
	\frac{v}{E_P} = \frac{\xi^4}{5\pi}, \qquad v = 245.65~\text{GeV} \quad (-0.23\,\%~\text{against}~246.22~\text{GeV}),
\end{equation}
dimensionless and without a fitted factor; the factor $10$ in $5\pi = (\pi/2)\cdot10$ is a motivated setting (Docs.~149, 383). The mass ratio $m_\mu/m_e = 207.8$ ($0.5\,\%$) from the previous section does not contain $v$ and is unaffected.

\textit{(Corrected on 30~Sept.~2026: The earlier derivation of the Higgs VEV has been removed -- the sections from \enquote{Derivation of Higgs VEV} to the final remark on fitting. It solved $m_e = y_e v$ for $v$ ($3.652\times10^{9}$ in natural units) and converted to GeV, where the conversion divided and multiplied by the same factor and thus gave $3.652\times10^{9}$~GeV instead of the stated $245.1$~GeV. The \enquote{compact formula} $v = \frac43\xi_0^{-1/2} = 115.5$ was converted with the unjustified divisor $10^{16}$ to \enquote{$141.0$~GeV} (correctly computed $1.41\times10^{5}$~GeV), and $K_{\text{quantum}} = 246.22/141.0 = 1.747$ was fixed from the measured value. The chain thus switched between different energy scales and yielded $v$ only through inserted or fitted values.)}

\section{List of Symbols Used}
\label{app:symbols_en}

\subsection{Fundamental Constants}
\adjustbox{max width=\textwidth}{%
\begin{tabular}{lll}
	\toprule
	\textbf{Symbol} & \textbf{Meaning} & \textbf{Value/Unit} \\
	\midrule
	$\xi$ & Geometric parameter & $\frac{4}{3} \times 10^{-4}$ (dimensionless) \\
	$c$ & Speed of light & $2.998 \times 10^8$ m/s \\
	$\hbar$ & Reduced Planck constant & $1.055 \times 10^{-34}$ J·s \\
	$G$ & Gravitational constant & $6.674 \times 10^{-11}$ m³/(kg·s²) \\
	$k_B$ & Boltzmann constant & $1.381 \times 10^{-23}$ J/K \\
	$e$ & Elementary charge & $1.602 \times 10^{-19}$ C \\
\end{tabular}%
}
\normalsize

\subsection{Coupling Constants}
\adjustbox{max width=\textwidth}{%
\begin{tabular}{lll}
	\toprule
	\textbf{Symbol} & \textbf{Meaning} & \textbf{Formula} \\
	\midrule
	$\alpha$ & Fine structure constant & $1/137.036$ (SI) \\
	$\alpha_{EM}$ & Electromagnetic coupling & $1$ (nat. units) \\
	$\alpha_S$ & Strong coupling & $\xi^{-1/3}$ \\
	$\alpha_W$ & Weak coupling & $\xi^{1/2}$ \\
	$\alpha_G$ & Gravitational coupling & $\xi^{2}$ \\
	$\varepsilon_T$ & T0 coupling parameter & $\xi \cdot E_0^2$ \\
	\bottomrule
\end{tabular}%
}
\normalsize

\subsection{Energy Scales and Masses}
\adjustbox{max width=\textwidth}{%
\begin{tabular}{lll}
	\toprule
	\textbf{Symbol} & \textbf{Meaning} & \textbf{Value/Formula} \\
	\midrule
	$E_P$ & Planck energy & $1.22 \times 10^{19}$ GeV \\
	$E_\xi$ & Characteristic energy & $1/\xi = 7500$ (nat. units) \\
	$E_0$ & Fundamental EM energy & $7.398$ MeV \\
	$v$ & Higgs VEV & $246.22$ GeV \\
	$m_h$ & Higgs mass & $125.25$ GeV \\
	$\Lambda_{QCD}$ & QCD scale & $\sim 200$ MeV \\
	$m_e$ & Electron mass & $0.511$ MeV \\
	$m_\mu$ & Muon mass & $105.66$ MeV \\
	$m_\tau$ & Tau mass & $1776.86$ MeV \\
	$m_u, m_d$ & Up, down quark masses & $2.16$, $4.67$ MeV \\
	$m_c, m_s$ & Charm, strange quark masses & $1.27$ GeV, $93.4$ MeV \\
	$m_t, m_b$ & Top, bottom quark masses & $172.76$ GeV, $4.18$ GeV \\
	$m_{\nu_e}, m_{\nu_\mu}, m_{\nu_\tau}$ & Neutrino masses & $< 2$ eV, $< 0.19$ MeV, $< 18.2$ MeV \\
	\bottomrule
\end{tabular}%
}
\normalsize

\subsection{Cosmological Parameters}
\adjustbox{max width=\textwidth}{%
\begin{tabular}{lll}
	\toprule
	\textbf{Symbol} & \textbf{Meaning} & \textbf{Value/Formula} \\
	\midrule
	$T_{CMB}$ & CMB temperature & $2.725$ K \\
	$z$ & Redshift & dimensionless \\
	$\Omega_\Lambda$ & Dark energy density & $0.6847$ ($\Lambda$CDM), $0$ (T0) \\
	$\Omega_{DM}$ & Dark matter density & $0.2607$ ($\Lambda$CDM), $0$ (T0) \\
	$\Omega_b$ & Baryon density & $0.0492$ ($\Lambda$CDM), $1$ (T0) \\
	$\Lambda$ & Cosmological constant & $(1.1 \pm 0.02) \times 10^{-52}$ m$^{-2}$ \\
	$\rho_\xi$ & $\xi$-field energy density & $E_\xi^4$ \\
	$\rho_{CMB}$ & CMB energy density & $4.64 \times 10^{-31}$ kg/m³ \\
	\bottomrule
\end{tabular}%
}
\normalsize

\subsection{Geometric and Derived Quantities}
\adjustbox{max width=\textwidth}{%
\begin{tabular}{lll}
	\toprule
	\textbf{Symbol} & \textbf{Meaning} & \textbf{Value/Formula} \\
	\midrule
	$D_f^{\text{eff}}$ & Effective fractal dimension (cumulative over RG running) & $2.973$ \\
	$\kappa_{mass}$ & Mass scaling exponent & $D_f^{\text{eff}}/2 = 1.487$ \\
	$\kappa_{grav}$ & Gravitational field parameter & $4.8 \times 10^{-11}$ m/s² \\
	$\lambda_h$ & Higgs self-coupling & $0.13$ \\
	$\theta_W$ & Weinberg angle & $\sin^2\theta_W = 0.2312$ \\
	$\theta_{QCD}$ & Strong CP phase & $< 10^{-10}$ (exp.), $\xi^2$ (T0) \\
	$\ell_P$ & Planck length & $1.616 \times 10^{-35}$ m \\
	$\lambda_C$ & Compton wavelength & $\hbar/(mc)$ \\
	$r_g$ & Gravitational radius & $2Gm$ \\
	$L_\xi$ & Characteristic length & $\xi$ (nat. units) \\
	\bottomrule
\end{tabular}%
}
\normalsize

\subsection{Mixing Matrices}
\adjustbox{max width=\textwidth}{%
\begin{tabular}{lll}
	\toprule
	\textbf{Symbol} & \textbf{Meaning} & \textbf{Typical Value} \\
	\midrule
	$V_{ij}$ & CKM matrix elements & see table \\
	$|V_{ud}|$ & CKM ud element & $0.97446$ \\
	$|V_{us}|$ & CKM us element (Cabibbo) & $0.22452$ \\
	$|V_{ub}|$ & CKM ub element & $0.00365$ \\
	$\delta_{CKM}$ & CKM CP phase & $1.20$ rad \\
	$\theta_{12}$ & PMNS solar angle & $33.44°$ \\
	$\theta_{23}$ & PMNS atmospheric & $49.2°$ \\
	$\theta_{13}$ & PMNS reactor angle & $8.57°$ \\
	$\delta_{CP}$ & PMNS CP phase & unknown \\
	\bottomrule
\end{tabular}%
}
\normalsize

\subsection{Other Symbols}
\adjustbox{max width=\textwidth}{%
\begin{tabular}{lll}
	\toprule
	\textbf{Symbol} & \textbf{Meaning} & \textbf{Context} \\
	\midrule
	$n, l, j$ & Quantum numbers & Particle classification \\
	$r_i$ & Rational coefficients & Yukawa couplings \\
	$p_i$ & Generation exponents & $3/2, 1, 2/3, ...$ \\
	$f(n,l,j)$ & Geometric function & Mass formula \\
	$\rho_{tet}$ & Tetrahedral packing density & $0.68$ \\
	$\gamma$ & Universal exponent & $1.01$ \\
	$\nu$ & Crystal symmetry factor & $0.63$ \\
	$\beta_T$ & Time field coupling & $1$ (nat. units) \\
	$y_i$ & Yukawa couplings & $r_i \cdot \xi^{p_i}$ \\
	$T(x,t)$ & Time field & T0 theory \\
	$E_{field}$ & Energy field & Universal field \\
	\bottomrule
\end{tabular}%
}
\normalsize	
\appendix

% ============================================================
% Addendum 7 September 2026 — Registered inconsistencies (K-041, R77)
% Source: Doc. 190-Archive-2, declared 25 August 2026 (K-041)
%         and Doc. 190 register R77
% Source document remains unchanged (append-only, cf. R50).
% ============================================================

\section*{Addendum: Internal Inconsistencies and Classification (7~September 2026)}

\textbf{Note:} The source document remains unchanged (append-only principle,
cf.\ Doc.~190, R50). The following entries document places where printed
formulas and table values in this document do not agree with the
document's own $\xi=(4/3)\cdot10^{-4}$. Resolution requires author
review of the source document.

\paragraph{K-041-a: CKM CP-phase $\delta_\mathrm{CKM}$ (declared 25~Aug.~2026).}
The printed formula $\delta_\mathrm{CKM}=\arcsin(2\sqrt{2}\cdot\sqrt{\xi}/3)$
yields $0.011$~rad, not the table value $1.20$~rad.
Candidate: $\arcsin(2\sqrt{2}/3)=1.231$~rad --- the factor $\sqrt{\xi}$
is likely a typesetting error. Whether this candidate is a corpus quantity
must be clarified by the author.

\paragraph{K-041-b: PMNS reactor angle $\theta_{13}$ (declared 25~Aug.~2026).}
The printed formula $\theta_{13}=\arcsin(\xi^{1/3})$ yields $2.93°$,
not the table value $8.57°$.
Candidate: argument $1/300=25\xi$ yields $8.59°$.
Whether $1/300$ is an independent corpus quantity must be clarified by the author.

\paragraph{K-041-c: CKM element $|V_{us}|$ (declared 25~Aug.~2026).}
The formula with Doc.~006 masses yields $|V_{us}|=0.2124$;
table value $0.22452$ (deviation $5.4\,\%$).
Possible cause: differing masses or $f_\mathrm{Cab}$ definition within this document.

\paragraph{R77: Classification of early coupling-constant formulas.}
The early coupling-constant formulas developed in this document are to be
classified as preliminary work, not as canonical FFGFT derivations.
Authoritative documents are Doc.~160 (lepton lifetime structure),
Doc.~318 (validity range of mass derivations), and the Galois series
(Docs.~336--348). The formulas in this document remain archived; numerical
use requires recourse to the successor documents.
